% =======================================================================
\section{Discussion}

\subsection{What Topological Persistence Adds for Practitioners}

The event-correlation result does not represent new tactical knowledge:
team shape breaks down under pressure and rebuilds during patient
possession. Its value is operational. Persistence provides a continuous,
automatically computed measure of tactical state that does not require
manual coding of match footage, and that can therefore be applied
consistently across a full match or a full season.

The detected loops correspond directly to tactical formations analysts
already describe qualitatively. A large tactical-scale loop, spanning
five to seven cluster centroids and persisting over many frames,
indicates that the attacking team is orbiting rather than penetrating
the defensive structure: the interior of the ring is a dead zone
controlled by the defence, consistent with the peripheral possession
ring described in Section~\ref{sec:intro}. The same class of large,
stable tactical loop also matches rest-defence arcs in positional play.
When deeper players hold a deliberate ring in a 3-2-5 or 2-3-5 shape,
the enclosed space is the counter-attack channel the rest defence is
designed to deny; persistence then measures how long that protective
geometry holds before a transition forces it to break. Rapid formation
and dissolution of smaller loops at the individual scale, by contrast,
corresponds to interlocking local cycles between a midfielder, winger,
and overlapping full-back, the kind of quick combination that breaks
down a peripheral ring by creating localised passing channels inside
the defensive block. A sharp decrease in individual-scale persistence
following a trigger event marks the closure of a pressing
encirclement: the ring tightens around the ball-carrier, and its
disappearance from the persistence diagram coincides with the moment
all exit angles are cut off. The size and duration of the preceding
loop quantifies how much space was enclosed before the press was
triggered.

The baseline comparison addresses a more specific concern: whether this
measure is redundant with geometric descriptors already in use, such as
team width or convex-hull area. The evidence indicates it is not.
Tactical-scale persistence correlates only weakly with the four
standard descriptors examined, and continues to explain variation in
team shape that those four measures do not account for jointly. This is
relevant to any decision about adding the measure to an existing
analytics pipeline, since it would not duplicate information already
captured. Whether the added information justifies the additional
computational and conceptual overhead is addressed by the next two
results.

\subsection{Tactical Independence and What It Suggests for Analysis}

The near-independence reported for home and away tactical persistence
is a useful negative finding rather than a non-finding. It establishes
a baseline: at frame-to-frame resolution, with no conditioning on match
context, one team's tactical structure carries almost no information
about the other's. Any future claim that tactical coupling exists
during specific phases of play, such as a sustained press, a passage of patient possession, or the moments around a set play, must demonstrate
an effect against this near-zero baseline, rather than against an
unstated assumption of independence. The present analysis cannot
distinguish between two explanations for the result: that no
meaningful coupling exists at any timescale, or that coupling exists
but only within specific tactical contexts that frame-by-frame
correlation across an entire match cannot detect. Testing the second
possibility would require conditioning the analysis on phase of play,
which the present ten-match sample is not large enough to support with
confidence.

\subsection{Why No Predictive Gain Yet, and What Would Change That}

The baseline comparison and the predictive-utility test are not in
tension, even though one shows tactical persistence adds information
and the other shows it does not improve prediction. The first is a
statement about variance within the present sample: after accounting
for the four geometric descriptors together, persistence still
explains some of what remains. The second is a stricter test, asking
whether that remaining explanatory power transfers to matches the model
has never seen. With only ten matches available to estimate how much a
single additional feature contributes, and with the geometric baseline
already predicting build-up phases reasonably well, the margin
available for any new feature to improve on is narrow. A genuine but
modest effect could fail to appear in held-out performance simply
because the sample is too small to estimate it reliably.

This distinction matters for how the result should be read. It is not
evidence that tactical-scale persistence lacks predictive value; it is
evidence that ten matches are not enough to detect whatever predictive
value it has, given how much of the predictive task the existing
descriptors already account for. Resolving this will require testing
the same comparison on a substantially larger sample, which is the
explicit aim of the companion full-season analysis.

\subsection{Limitations}

The present findings are drawn from ten matches in a single
competition, the A-League, and that scope limits how far the results
should be generalised. Tactical systems, squad quality, and playing
conditions vary between leagues, and a measure validated here may
behave differently in competitions with markedly different tempos or
spacing conventions. The companion full-season analysis, drawn from a
single Championship season, will extend the evidence base considerably
but will not by itself resolve cross-competition generalisability; that
would require matches from multiple leagues, which is not yet part of
either dataset.

All tracking data analysed here are broadcast-derived, at 10~Hz. Clubs
using higher-frequency optical tracking systems, typically around
25~Hz, may find that persistence values computed on their own data do
not match the magnitudes reported here, even if the qualitative
patterns, including events tracking persistence changes and weak
coupling with geometric descriptors, hold. Direct comparison across tracking
technologies has not been tested and would be needed before treating
the specific numerical thresholds reported here as transferable to a
different data source.

The predictive-utility test used a single model family, L2-regularised
logistic regression, and a single way of combining tactical persistence
with the geometric baseline. Other model architectures or feature
combinations might extract value from tactical persistence that this
particular specification does not, although the consistency of the null result across three different persistence summaries, total, count, and maximum, makes a large undetected effect from model
choice alone less likely.

\subsection{Future Directions in Other Sports}

The method validated in the companion paper is not specific to
football, and the questions asked here translate readily to other
invasion sports with comparable spatial tracking data. Rugby union and
rugby league involve similar individual-tactical-team scale structure,
though pitch dimensions and player counts differ enough that the
validated cutoff distances reported here would need re-deriving rather
than reused directly. Basketball presents a different scale problem:
the court is small enough, and player counts low enough, that the
individual and tactical scales identified in football may partially
collapse into one, which would need testing rather than assuming. In
each case, the same two-step validation used here, establishing stable scale regimes before asking what they reveal about play, would
need repeating on sport-specific data before any of the present
findings could be assumed to transfer. We see this as a natural next
step rather than a claim made by the present work.
